\chapter{Feature Engineering at Scale}
\label{ch:feature_engineering}

\section*{Learning Objectives}

By the end of this chapter, you will be able to:

\begin{itemize}
    \item Design systematic feature engineering pipelines for production ML systems
    \item Extract temporal features from timestamp data including cyclical and holiday features
    \item Implement scalable aggregation and rolling window features for time series and grouped data
    \item Create interaction and polynomial features while managing dimensionality
    \item Apply advanced categorical encoding strategies including target encoding and embeddings
    \item Implement feature selection techniques to reduce dimensionality and improve model performance
    \item Build automated feature engineering frameworks for reproducibility and maintainability
\end{itemize}

\section{Introduction}

Feature engineering is often cited as the most impactful aspect of applied machine learning. While model selection and hyperparameter tuning provide incremental improvements, good features can transform an intractable problem into a solvable one.

In production ML systems, feature engineering must be:

\begin{itemize}
    \item \textbf{Scalable:} Handle large datasets efficiently
    \item \textbf{Reproducible:} Consistent across training and serving
    \item \textbf{Maintainable:} Clear code and documentation
    \item \textbf{Robust:} Handle missing values and edge cases
    \item \textbf{Monitored:} Track feature distributions and drift
\end{itemize}

\subsection{Feature Engineering Principles}

\begin{table}[h]
\centering
\caption{Feature Engineering Best Practices}
\label{tab:feature_engineering_principles}
\begin{tabular}{|l|p{10cm}|}
\hline
\textbf{Principle} & \textbf{Description} \\
\hline
Domain Knowledge & Leverage business understanding to create meaningful features \\
Train-Test Consistency & Ensure features computed identically during training and serving \\
Avoid Data Leakage & Never use future information or target variables \\
Handle Missing Values & Explicit strategy for nulls and edge cases \\
Feature Documentation & Clear naming and metadata for all features \\
Computational Efficiency & Consider computation cost for real-time serving \\
\hline
\end{tabular}
\end{table}

\subsection{Feature Engineering Pipeline Architecture}

\begin{figure}[h]
\centering
\begin{tikzpicture}[
    node distance=1.5cm,
    box/.style={rectangle, draw, text width=3cm, align=center, minimum height=1cm},
    arrow/.style={->, >=stealth, thick}
]
    \node[box] (raw) {Raw Data};
    \node[box, right=of raw] (clean) {Data Cleaning};
    \node[box, right=of clean] (temporal) {Temporal Features};
    \node[box, below=of temporal] (agg) {Aggregations};
    \node[box, left=of agg] (encode) {Encoding};
    \node[box, left=of encode] (select) {Feature Selection};
    \node[box, below=of encode] (output) {Feature Matrix};

    \draw[arrow] (raw) -- (clean);
    \draw[arrow] (clean) -- (temporal);
    \draw[arrow] (temporal) -- (agg);
    \draw[arrow] (agg) -- (encode);
    \draw[arrow] (encode) -- (select);
    \draw[arrow] (select) -- (output);
\end{tikzpicture}
\caption{Feature Engineering Pipeline Flow}
\label{fig:feature_pipeline}
\end{figure}

\section{Feature Engineering Fundamentals}

\subsection{Comprehensive Feature Engineering Framework}

\begin{lstlisting}[style=python, caption={Production Feature Engineering Framework}]
from typing import Dict, List, Optional, Union, Callable, Any, Tuple
from dataclasses import dataclass, field
import pandas as pd
import numpy as np
from sklearn.base import BaseEstimator, TransformerMixin
from sklearn.preprocessing import StandardScaler, RobustScaler
import logging
from pathlib import Path
import json
import hashlib
from datetime import datetime

@dataclass
class FeatureConfig:
    """Configuration for feature engineering pipeline."""
    name: str
    feature_type: str  # 'numerical', 'categorical', 'temporal', 'text'
    transformations: List[str] = field(default_factory=list)
    required: bool = True
    default_value: Any = None
    scaling_strategy: Optional[str] = None  # 'standard', 'robust', 'minmax'
    encoding_strategy: Optional[str] = None  # 'onehot', 'label', 'target', 'frequency'

@dataclass
class FeatureMetadata:
    """Metadata for tracking feature provenance."""
    name: str
    source_columns: List[str]
    transformation_logic: str
    created_at: datetime
    version: str
    statistics: Dict[str, float] = field(default_factory=dict)

class FeatureEngineer(BaseEstimator, TransformerMixin):
    """
    Comprehensive feature engineering framework.

    Features:
    - Modular transformation pipeline
    - Feature versioning and metadata tracking
    - Train/test consistency enforcement
    - Missing value handling
    - Feature validation
    """

    def __init__(
        self,
        feature_configs: Optional[List[FeatureConfig]] = None,
        handle_missing: str = 'median',  # 'median', 'mean', 'mode', 'drop', 'constant'
        missing_indicator: bool = True,
        validate_features: bool = True,
        track_metadata: bool = True
    ):
        """
        Initialize feature engineer.

        Args:
            feature_configs: List of feature configurations
            handle_missing: Strategy for handling missing values
            missing_indicator: Add indicator features for missing values
            validate_features: Enable feature validation
            track_metadata: Enable metadata tracking
        """
        self.feature_configs = feature_configs or []
        self.handle_missing = handle_missing
        self.missing_indicator = missing_indicator
        self.validate_features = validate_features
        self.track_metadata = track_metadata

        # State
        self.feature_metadata_: Dict[str, FeatureMetadata] = {}
        self.missing_value_fills_: Dict[str, Any] = {}
        self.fitted_ = False

        self.logger = logging.getLogger(__name__)

    def fit(self, X: pd.DataFrame, y: Optional[pd.Series] = None) -> 'FeatureEngineer':
        """
        Fit feature engineering pipeline.

        Args:
            X: Input data
            y: Target variable (optional, needed for target encoding)

        Returns:
            Self
        """
        # Compute missing value fills
        self._compute_missing_fills(X)

        # Store target for target encoding if provided
        self.y_fit_ = y

        self.fitted_ = True
        self.logger.info(f"Feature engineer fitted on {len(X)} samples")

        return self

    def transform(self, X: pd.DataFrame) -> pd.DataFrame:
        """
        Transform data using fitted feature engineering pipeline.

        Args:
            X: Input data

        Returns:
            Transformed data
        """
        if not self.fitted_:
            raise ValueError("FeatureEngineer must be fitted before transform")

        X_transformed = X.copy()

        # Handle missing values
        X_transformed = self._handle_missing_values(X_transformed)

        # Validate features if enabled
        if self.validate_features:
            self._validate_features(X_transformed)

        return X_transformed

    def fit_transform(self, X: pd.DataFrame, y: Optional[pd.Series] = None) -> pd.DataFrame:
        """Fit and transform in one step."""
        return self.fit(X, y).transform(X)

    def _compute_missing_fills(self, X: pd.DataFrame):
        """Compute values for filling missing data."""
        for col in X.columns:
            if X[col].isna().any():
                if self.handle_missing == 'median' and pd.api.types.is_numeric_dtype(X[col]):
                    self.missing_value_fills_[col] = X[col].median()
                elif self.handle_missing == 'mean' and pd.api.types.is_numeric_dtype(X[col]):
                    self.missing_value_fills_[col] = X[col].mean()
                elif self.handle_missing == 'mode':
                    self.missing_value_fills_[col] = X[col].mode()[0] if not X[col].mode().empty else None
                elif self.handle_missing == 'constant':
                    self.missing_value_fills_[col] = 0 if pd.api.types.is_numeric_dtype(X[col]) else 'MISSING'
                else:
                    self.missing_value_fills_[col] = None

    def _handle_missing_values(self, X: pd.DataFrame) -> pd.DataFrame:
        """Handle missing values according to strategy."""
        X_filled = X.copy()

        for col, fill_value in self.missing_value_fills_.items():
            if col in X_filled.columns:
                # Add missing indicator if enabled
                if self.missing_indicator and X_filled[col].isna().any():
                    X_filled[f'{col}_is_missing'] = X_filled[col].isna().astype(int)

                # Fill missing values
                if fill_value is not None:
                    X_filled[col].fillna(fill_value, inplace=True)

        return X_filled

    def _validate_features(self, X: pd.DataFrame):
        """Validate feature quality and distributions."""
        for col in X.columns:
            # Check for infinite values
            if pd.api.types.is_numeric_dtype(X[col]):
                if np.isinf(X[col]).any():
                    self.logger.warning(f"Feature {col} contains infinite values")

            # Check for high cardinality in categorical features
            if pd.api.types.is_object_dtype(X[col]) or pd.api.types.is_categorical_dtype(X[col]):
                n_unique = X[col].nunique()
                if n_unique > 1000:
                    self.logger.warning(
                        f"Feature {col} has high cardinality: {n_unique} unique values"
                    )

            # Check for constant features
            if X[col].nunique() == 1:
                self.logger.warning(f"Feature {col} is constant")

    def add_feature_metadata(
        self,
        name: str,
        source_columns: List[str],
        transformation_logic: str,
        statistics: Optional[Dict[str, float]] = None
    ):
        """Add metadata for a feature."""
        if self.track_metadata:
            self.feature_metadata_[name] = FeatureMetadata(
                name=name,
                source_columns=source_columns,
                transformation_logic=transformation_logic,
                created_at=datetime.now(),
                version="1.0",
                statistics=statistics or {}
            )

    def get_feature_metadata(self) -> Dict[str, Dict[str, Any]]:
        """Get all feature metadata."""
        return {
            name: {
                'source_columns': meta.source_columns,
                'transformation_logic': meta.transformation_logic,
                'created_at': meta.created_at.isoformat(),
                'version': meta.version,
                'statistics': meta.statistics
            }
            for name, meta in self.feature_metadata_.items()
        }

    def save_state(self, path: str):
        """Save feature engineer state for reproducibility."""
        state = {
            'missing_value_fills': self.missing_value_fills_,
            'feature_metadata': self.get_feature_metadata(),
            'fitted': self.fitted_
        }

        with open(path, 'w') as f:
            json.dump(state, f, indent=2, default=str)

        self.logger.info(f"Feature engineer state saved to {path}")

    def load_state(self, path: str):
        """Load feature engineer state."""
        with open(path, 'r') as f:
            state = json.load(f)

        self.missing_value_fills_ = state['missing_value_fills']
        self.fitted_ = state['fitted']

        self.logger.info(f"Feature engineer state loaded from {path}")

class FeatureStore:
    """
    Simple feature store for managing feature definitions and values.

    In production, consider using dedicated feature stores like:
    - Feast
    - Tecton
    - AWS Feature Store
    - Databricks Feature Store
    """

    def __init__(self, storage_path: str = "./feature_store"):
        """
        Initialize feature store.

        Args:
            storage_path: Path for storing feature definitions and values
        """
        self.storage_path = Path(storage_path)
        self.storage_path.mkdir(parents=True, exist_ok=True)
        self.logger = logging.getLogger(__name__)

        # Feature registry
        self.feature_registry: Dict[str, FeatureConfig] = {}

    def register_feature(self, config: FeatureConfig):
        """Register a feature definition."""
        self.feature_registry[config.name] = config
        self.logger.info(f"Registered feature: {config.name}")

        # Save to disk
        self._save_registry()

    def get_feature(self, name: str) -> Optional[FeatureConfig]:
        """Get feature configuration by name."""
        return self.feature_registry.get(name)

    def list_features(self) -> List[str]:
        """List all registered features."""
        return list(self.feature_registry.keys())

    def _save_registry(self):
        """Save feature registry to disk."""
        registry_path = self.storage_path / "feature_registry.json"
        registry_data = {
            name: {
                'name': config.name,
                'feature_type': config.feature_type,
                'transformations': config.transformations,
                'required': config.required,
                'default_value': config.default_value,
                'scaling_strategy': config.scaling_strategy,
                'encoding_strategy': config.encoding_strategy
            }
            for name, config in self.feature_registry.items()
        }

        with open(registry_path, 'w') as f:
            json.dump(registry_data, f, indent=2)

    def compute_features(
        self,
        entity_df: pd.DataFrame,
        feature_names: List[str],
        timestamp_column: Optional[str] = None
    ) -> pd.DataFrame:
        """
        Compute features for entities.

        Args:
            entity_df: DataFrame with entity IDs
            feature_names: List of features to compute
            timestamp_column: Column for point-in-time correctness

        Returns:
            DataFrame with computed features
        """
        # In production, this would fetch precomputed features
        # For now, return placeholder
        self.logger.info(f"Computing {len(feature_names)} features for {len(entity_df)} entities")
        return entity_df

# Example usage
if __name__ == "__main__":
    # Generate sample data
    np.random.seed(42)
    df = pd.DataFrame({
        'customer_id': range(1000),
        'age': np.random.randint(18, 80, 1000),
        'income': np.random.normal(50000, 20000, 1000),
        'purchase_amount': np.random.exponential(100, 1000),
        'category': np.random.choice(['A', 'B', 'C', 'D'], 1000),
        'timestamp': pd.date_range('2023-01-01', periods=1000, freq='H')
    })

    # Add missing values
    df.loc[df.sample(100).index, 'income'] = np.nan
    df.loc[df.sample(50).index, 'category'] = np.nan

    # Initialize feature engineer
    engineer = FeatureEngineer(
        handle_missing='median',
        missing_indicator=True,
        validate_features=True
    )

    # Fit and transform
    df_transformed = engineer.fit_transform(df)
    print(f"\nTransformed shape: {df_transformed.shape}")
    print(f"Missing value fills: {engineer.missing_value_fills_}")

    # Save state
    engineer.save_state("feature_engineer_state.json")
\end{lstlisting}

\section{Temporal and Time-based Features}

Temporal features are crucial for time series forecasting, user behavior modeling, and event prediction. Proper temporal feature extraction can capture seasonality, trends, and cyclical patterns.

\subsection{Temporal Feature Extraction}

\begin{lstlisting}[style=python, caption={Comprehensive Temporal Feature Engineering}]
from typing import List, Optional
import pandas as pd
import numpy as np
from datetime import datetime, timedelta

class TemporalFeatureExtractor:
    """
    Extract comprehensive temporal features from datetime columns.

    Features extracted:
    - Basic components (year, month, day, hour, etc.)
    - Cyclical encodings (sin/cos transformations)
    - Business calendar features (holidays, weekends)
    - Relative time features (time since event)
    - Temporal aggregations
    """

    def __init__(
        self,
        timestamp_column: str,
        country_holidays: str = 'US',  # Country code for holidays
        extract_cyclical: bool = True,
        extract_business_features: bool = True
    ):
        """
        Initialize temporal feature extractor.

        Args:
            timestamp_column: Name of timestamp column
            country_holidays: Country code for holiday calendar
            extract_cyclical: Extract cyclical (sin/cos) features
            extract_business_features: Extract business calendar features
        """
        self.timestamp_column = timestamp_column
        self.country_holidays = country_holidays
        self.extract_cyclical = extract_cyclical
        self.extract_business_features = extract_business_features

        # Try to load holidays library
        try:
            import holidays
            self.holidays = holidays.country_holidays(country_holidays)
        except ImportError:
            self.holidays = None
            if extract_business_features:
                logging.warning("holidays library not installed, business features will be limited")

    def extract_basic_features(self, df: pd.DataFrame) -> pd.DataFrame:
        """
        Extract basic temporal components.

        Args:
            df: DataFrame with timestamp column

        Returns:
            DataFrame with added temporal features
        """
        df = df.copy()
        ts = pd.to_datetime(df[self.timestamp_column])

        # Basic components
        df['year'] = ts.dt.year
        df['month'] = ts.dt.month
        df['day'] = ts.dt.day
        df['dayofweek'] = ts.dt.dayofweek  # 0 = Monday
        df['dayofyear'] = ts.dt.dayofyear
        df['quarter'] = ts.dt.quarter
        df['weekofyear'] = ts.dt.isocalendar().week
        df['hour'] = ts.dt.hour
        df['minute'] = ts.dt.minute

        # Boolean indicators
        df['is_weekend'] = (ts.dt.dayofweek >= 5).astype(int)
        df['is_month_start'] = ts.dt.is_month_start.astype(int)
        df['is_month_end'] = ts.dt.is_month_end.astype(int)
        df['is_quarter_start'] = ts.dt.is_quarter_start.astype(int)
        df['is_quarter_end'] = ts.dt.is_quarter_end.astype(int)

        return df

    def extract_cyclical_features(self, df: pd.DataFrame) -> pd.DataFrame:
        """
        Extract cyclical features using sin/cos transformations.

        Cyclical encoding ensures smooth transitions (e.g., 11pm -> 12am).

        Args:
            df: DataFrame with timestamp column

        Returns:
            DataFrame with cyclical features
        """
        if not self.extract_cyclical:
            return df

        df = df.copy()
        ts = pd.to_datetime(df[self.timestamp_column])

        # Month cyclical (12 months)
        df['month_sin'] = np.sin(2 * np.pi * ts.dt.month / 12)
        df['month_cos'] = np.cos(2 * np.pi * ts.dt.month / 12)

        # Day of week cyclical (7 days)
        df['dayofweek_sin'] = np.sin(2 * np.pi * ts.dt.dayofweek / 7)
        df['dayofweek_cos'] = np.cos(2 * np.pi * ts.dt.dayofweek / 7)

        # Day of month cyclical (assuming 31 days)
        df['day_sin'] = np.sin(2 * np.pi * ts.dt.day / 31)
        df['day_cos'] = np.cos(2 * np.pi * ts.dt.day / 31)

        # Hour cyclical (24 hours)
        df['hour_sin'] = np.sin(2 * np.pi * ts.dt.hour / 24)
        df['hour_cos'] = np.cos(2 * np.pi * ts.dt.hour / 24)

        # Day of year cyclical (365 days)
        df['dayofyear_sin'] = np.sin(2 * np.pi * ts.dt.dayofyear / 365)
        df['dayofyear_cos'] = np.cos(2 * np.pi * ts.dt.dayofyear / 365)

        return df

    def extract_business_features(self, df: pd.DataFrame) -> pd.DataFrame:
        """
        Extract business calendar features.

        Args:
            df: DataFrame with timestamp column

        Returns:
            DataFrame with business features
        """
        if not self.extract_business_features:
            return df

        df = df.copy()
        ts = pd.to_datetime(df[self.timestamp_column])

        # Business day indicator
        df['is_business_day'] = ts.dt.dayofweek < 5  # Monday-Friday

        # Holiday indicator
        if self.holidays is not None:
            df['is_holiday'] = ts.dt.date.apply(lambda d: d in self.holidays).astype(int)

            # Days until/since holiday
            df['days_to_next_holiday'] = ts.apply(self._days_to_next_holiday)
            df['days_since_last_holiday'] = ts.apply(self._days_since_last_holiday)
        else:
            df['is_holiday'] = 0
            df['days_to_next_holiday'] = -1
            df['days_since_last_holiday'] = -1

        # Time of day categories
        hour = ts.dt.hour
        df['time_of_day'] = 'night'  # 0-6
        df.loc[hour >= 6, 'time_of_day'] = 'morning'  # 6-12
        df.loc[hour >= 12, 'time_of_day'] = 'afternoon'  # 12-18
        df.loc[hour >= 18, 'time_of_day'] = 'evening'  # 18-24

        # Part of month
        day = ts.dt.day
        df['part_of_month'] = 'beginning'  # 1-10
        df.loc[day > 10, 'part_of_month'] = 'middle'  # 11-20
        df.loc[day > 20, 'part_of_month'] = 'end'  # 21-31

        return df

    def _days_to_next_holiday(self, dt: pd.Timestamp) -> int:
        """Calculate days to next holiday."""
        if self.holidays is None:
            return -1

        current_date = dt.date()
        for i in range(1, 366):  # Look ahead up to 1 year
            future_date = current_date + timedelta(days=i)
            if future_date in self.holidays:
                return i
        return 365

    def _days_since_last_holiday(self, dt: pd.Timestamp) -> int:
        """Calculate days since last holiday."""
        if self.holidays is None:
            return -1

        current_date = dt.date()
        for i in range(1, 366):  # Look back up to 1 year
            past_date = current_date - timedelta(days=i)
            if past_date in self.holidays:
                return i
        return 365

    def extract_relative_features(
        self,
        df: pd.DataFrame,
        reference_date: Optional[datetime] = None,
        event_column: Optional[str] = None
    ) -> pd.DataFrame:
        """
        Extract relative time features.

        Args:
            df: DataFrame with timestamp column
            reference_date: Reference date for time calculations
            event_column: Column containing event timestamps

        Returns:
            DataFrame with relative time features
        """
        df = df.copy()
        ts = pd.to_datetime(df[self.timestamp_column])

        # Time since reference date
        if reference_date:
            ref_date = pd.to_datetime(reference_date)
            df['days_since_reference'] = (ts - ref_date).dt.days
            df['hours_since_reference'] = (ts - ref_date).dt.total_seconds() / 3600

        # Time since event
        if event_column and event_column in df.columns:
            event_ts = pd.to_datetime(df[event_column])
            df['days_since_event'] = (ts - event_ts).dt.days
            df['hours_since_event'] = (ts - event_ts).dt.total_seconds() / 3600

        # Time to end of period
        df['days_until_month_end'] = (ts + pd.offsets.MonthEnd(0) - ts).dt.days
        df['days_until_quarter_end'] = (ts + pd.offsets.QuarterEnd(0) - ts).dt.days
        df['days_until_year_end'] = (ts + pd.offsets.YearEnd(0) - ts).dt.days

        return df

    def extract_all_features(self, df: pd.DataFrame, **kwargs) -> pd.DataFrame:
        """
        Extract all temporal features.

        Args:
            df: DataFrame with timestamp column
            **kwargs: Additional arguments for specific extractors

        Returns:
            DataFrame with all temporal features
        """
        df = self.extract_basic_features(df)
        df = self.extract_cyclical_features(df)
        df = self.extract_business_features(df)
        df = self.extract_relative_features(df, **kwargs)

        return df

# Example usage
if __name__ == "__main__":
    # Generate sample data
    df = pd.DataFrame({
        'timestamp': pd.date_range('2023-01-01', periods=1000, freq='H'),
        'value': np.random.randn(1000)
    })

    # Extract temporal features
    extractor = TemporalFeatureExtractor(
        timestamp_column='timestamp',
        country_holidays='US',
        extract_cyclical=True,
        extract_business_features=True
    )

    df_features = extractor.extract_all_features(df)
    print(f"\nOriginal columns: {df.shape[1]}")
    print(f"After feature extraction: {df_features.shape[1]}")
    print(f"\nNew features: {[col for col in df_features.columns if col not in df.columns]}")
\end{lstlisting}

\section{Aggregation and Rolling Features}

Aggregation and rolling window features are essential for capturing historical patterns and trends in time series and grouped data.

\subsection{Scalable Aggregation Features}

\begin{lstlisting}[style=python, caption={Scalable Aggregation and Rolling Window Features}]
from typing import List, Dict, Optional, Union, Callable
import pandas as pd
import numpy as np
from functools import partial

class AggregationFeatureGenerator:
    """
    Generate aggregation and rolling window features at scale.

    Features:
    - Group-by aggregations
    - Rolling window statistics
    - Expanding window features
    - Lag features
    - Custom aggregation functions
    """

    def __init__(
        self,
        group_columns: Optional[List[str]] = None,
        timestamp_column: Optional[str] = None
    ):
        """
        Initialize aggregation feature generator.

        Args:
            group_columns: Columns to group by
            timestamp_column: Column for temporal operations
        """
        self.group_columns = group_columns or []
        self.timestamp_column = timestamp_column

    def create_group_aggregations(
        self,
        df: pd.DataFrame,
        value_columns: List[str],
        agg_functions: List[Union[str, Callable]] = ['mean', 'std', 'min', 'max', 'count']
    ) -> pd.DataFrame:
        """
        Create group-by aggregation features.

        Args:
            df: Input DataFrame
            value_columns: Columns to aggregate
            agg_functions: Aggregation functions to apply

        Returns:
            DataFrame with aggregation features
        """
        if not self.group_columns:
            raise ValueError("group_columns must be specified")

        df = df.copy()
        agg_dict = {col: agg_functions for col in value_columns}

        # Compute aggregations
        agg_df = df.groupby(self.group_columns).agg(agg_dict)

        # Flatten column names
        agg_df.columns = ['_'.join(col).strip() for col in agg_df.columns.values]

        # Merge back to original DataFrame
        df = df.merge(agg_df, on=self.group_columns, how='left')

        return df

    def create_rolling_features(
        self,
        df: pd.DataFrame,
        value_column: str,
        windows: List[int] = [7, 14, 30],
        functions: List[str] = ['mean', 'std', 'min', 'max']
    ) -> pd.DataFrame:
        """
        Create rolling window features.

        Args:
            df: Input DataFrame
            value_column: Column to compute rolling features on
            windows: List of window sizes
            functions: List of aggregation functions

        Returns:
            DataFrame with rolling features
        """
        if self.timestamp_column is None:
            raise ValueError("timestamp_column must be specified for rolling features")

        df = df.copy()
        df = df.sort_values(self.timestamp_column)

        for window in windows:
            for func in functions:
                feature_name = f'{value_column}_rolling_{window}_{func}'

                if self.group_columns:
                    # Rolling within groups
                    df[feature_name] = df.groupby(self.group_columns)[value_column].transform(
                        lambda x: x.rolling(window=window, min_periods=1).agg(func)
                    )
                else:
                    # Global rolling
                    df[feature_name] = df[value_column].rolling(
                        window=window, min_periods=1
                    ).agg(func)

        return df

    def create_lag_features(
        self,
        df: pd.DataFrame,
        value_column: str,
        lags: List[int] = [1, 7, 14, 30]
    ) -> pd.DataFrame:
        """
        Create lag features.

        Args:
            df: Input DataFrame
            value_column: Column to create lags for
            lags: List of lag periods

        Returns:
            DataFrame with lag features
        """
        if self.timestamp_column is None:
            raise ValueError("timestamp_column must be specified for lag features")

        df = df.copy()
        df = df.sort_values(self.timestamp_column)

        for lag in lags:
            feature_name = f'{value_column}_lag_{lag}'

            if self.group_columns:
                # Lag within groups
                df[feature_name] = df.groupby(self.group_columns)[value_column].shift(lag)
            else:
                # Global lag
                df[feature_name] = df[value_column].shift(lag)

        return df

    def create_expanding_features(
        self,
        df: pd.DataFrame,
        value_column: str,
        functions: List[str] = ['mean', 'std', 'min', 'max', 'count']
    ) -> pd.DataFrame:
        """
        Create expanding window features (cumulative statistics).

        Args:
            df: Input DataFrame
            value_column: Column to compute expanding features on
            functions: List of aggregation functions

        Returns:
            DataFrame with expanding features
        """
        if self.timestamp_column is None:
            raise ValueError("timestamp_column must be specified for expanding features")

        df = df.copy()
        df = df.sort_values(self.timestamp_column)

        for func in functions:
            feature_name = f'{value_column}_expanding_{func}'

            if self.group_columns:
                # Expanding within groups
                df[feature_name] = df.groupby(self.group_columns)[value_column].transform(
                    lambda x: x.expanding(min_periods=1).agg(func)
                )
            else:
                # Global expanding
                df[feature_name] = df[value_column].expanding(min_periods=1).agg(func)

        return df

    def create_diff_features(
        self,
        df: pd.DataFrame,
        value_column: str,
        periods: List[int] = [1, 7]
    ) -> pd.DataFrame:
        """
        Create difference features (change from previous period).

        Args:
            df: Input DataFrame
            value_column: Column to compute differences for
            periods: List of periods for differences

        Returns:
            DataFrame with difference features
        """
        if self.timestamp_column is None:
            raise ValueError("timestamp_column must be specified for diff features")

        df = df.copy()
        df = df.sort_values(self.timestamp_column)

        for period in periods:
            feature_name = f'{value_column}_diff_{period}'

            if self.group_columns:
                # Diff within groups
                df[feature_name] = df.groupby(self.group_columns)[value_column].diff(period)
            else:
                # Global diff
                df[feature_name] = df[value_column].diff(period)

            # Percentage change
            pct_change_name = f'{value_column}_pct_change_{period}'
            if self.group_columns:
                df[pct_change_name] = df.groupby(self.group_columns)[value_column].pct_change(period)
            else:
                df[pct_change_name] = df[value_column].pct_change(period)

        return df

    def create_time_based_aggregations(
        self,
        df: pd.DataFrame,
        value_column: str,
        time_windows: List[str] = ['1D', '7D', '30D'],
        functions: List[str] = ['mean', 'std', 'sum']
    ) -> pd.DataFrame:
        """
        Create time-based rolling aggregations.

        Args:
            df: Input DataFrame
            value_column: Column to aggregate
            time_windows: List of time windows (pandas offset strings)
            functions: List of aggregation functions

        Returns:
            DataFrame with time-based aggregations
        """
        if self.timestamp_column is None:
            raise ValueError("timestamp_column must be specified")

        df = df.copy()
        df = df.sort_values(self.timestamp_column)

        # Set timestamp as index for time-based rolling
        df_indexed = df.set_index(self.timestamp_column)

        for window in time_windows:
            for func in functions:
                feature_name = f'{value_column}_rolling_{window}_{func}'

                if self.group_columns:
                    # Time-based rolling within groups
                    df[feature_name] = df_indexed.groupby(self.group_columns)[value_column].transform(
                        lambda x: x.rolling(window=window, min_periods=1).agg(func)
                    )
                else:
                    # Global time-based rolling
                    df[feature_name] = df_indexed[value_column].rolling(
                        window=window, min_periods=1
                    ).agg(func).values

        return df

    def create_all_aggregation_features(
        self,
        df: pd.DataFrame,
        value_columns: List[str],
        rolling_windows: List[int] = [7, 14, 30],
        lags: List[int] = [1, 7, 14],
        time_windows: List[str] = ['7D', '30D']
    ) -> pd.DataFrame:
        """
        Create comprehensive set of aggregation features.

        Args:
            df: Input DataFrame
            value_columns: Columns to create features for
            rolling_windows: Rolling window sizes
            lags: Lag periods
            time_windows: Time-based windows

        Returns:
            DataFrame with all aggregation features
        """
        df_result = df.copy()

        for value_col in value_columns:
            # Rolling features
            if rolling_windows:
                df_result = self.create_rolling_features(
                    df_result, value_col, windows=rolling_windows
                )

            # Lag features
            if lags:
                df_result = self.create_lag_features(df_result, value_col, lags=lags)

            # Expanding features
            df_result = self.create_expanding_features(df_result, value_col)

            # Diff features
            df_result = self.create_diff_features(df_result, value_col)

            # Time-based aggregations
            if time_windows and self.timestamp_column:
                df_result = self.create_time_based_aggregations(
                    df_result, value_col, time_windows=time_windows
                )

        return df_result

# Example usage
if __name__ == "__main__":
    # Generate sample time series data
    np.random.seed(42)
    dates = pd.date_range('2023-01-01', periods=365, freq='D')
    customers = ['A', 'B', 'C']

    df = pd.DataFrame({
        'date': np.repeat(dates, len(customers)),
        'customer': np.tile(customers, len(dates)),
        'sales': np.random.exponential(100, len(dates) * len(customers)),
        'quantity': np.random.poisson(10, len(dates) * len(customers))
    })

    # Initialize generator
    generator = AggregationFeatureGenerator(
        group_columns=['customer'],
        timestamp_column='date'
    )

    # Create all features
    df_features = generator.create_all_aggregation_features(
        df,
        value_columns=['sales', 'quantity'],
        rolling_windows=[7, 14, 30],
        lags=[1, 7, 14, 30],
        time_windows=['7D', '30D']
    )

    print(f"\nOriginal shape: {df.shape}")
    print(f"With features: {df_features.shape}")
    print(f"\nNew features ({df_features.shape[1] - df.shape[1]} total):")
    new_features = [col for col in df_features.columns if col not in df.columns]
    for feat in new_features[:10]:  # Show first 10
        print(f"  - {feat}")
\end{lstlisting}

\section{Interaction and Polynomial Features}

Interaction and polynomial features can capture non-linear relationships and feature interactions, but must be managed carefully to avoid dimensionality explosion.

\subsection{Feature Interactions}

\begin{lstlisting}[style=python, caption={Intelligent Feature Interaction Generation}]
from typing import List, Tuple, Optional, Dict
import pandas as pd
import numpy as np
from sklearn.preprocessing import PolynomialFeatures
from itertools import combinations
import logging

class InteractionFeatureGenerator:
    """
    Generate interaction and polynomial features with dimensionality management.

    Features:
    - Pairwise interactions
    - Polynomial features
    - Ratio features
    - Product features
    - Feature selection based on importance
    """

    def __init__(
        self,
        max_interactions: int = 100,
        interaction_degree: int = 2,
        include_polynomial: bool = True,
        polynomial_degree: int = 2,
        include_bias: bool = False
    ):
        """
        Initialize interaction feature generator.

        Args:
            max_interactions: Maximum number of interaction features to create
            interaction_degree: Degree of interactions (2 = pairwise)
            include_polynomial: Include polynomial features
            polynomial_degree: Degree of polynomial features
            include_bias: Include bias term in polynomial features
        """
        self.max_interactions = max_interactions
        self.interaction_degree = interaction_degree
        self.include_polynomial = include_polynomial
        self.polynomial_degree = polynomial_degree
        self.include_bias = include_bias

        self.logger = logging.getLogger(__name__)
        self.selected_interactions_: Optional[List[Tuple[str, ...]]] = None

    def create_pairwise_interactions(
        self,
        df: pd.DataFrame,
        numerical_columns: List[str],
        interaction_type: str = 'multiply'  # 'multiply', 'add', 'subtract', 'divide'
    ) -> pd.DataFrame:
        """
        Create pairwise interaction features.

        Args:
            df: Input DataFrame
            numerical_columns: Numerical columns to create interactions for
            interaction_type: Type of interaction

        Returns:
            DataFrame with interaction features
        """
        df = df.copy()

        # Generate all pairwise combinations
        pairs = list(combinations(numerical_columns, 2))

        # Limit number of interactions
        if len(pairs) > self.max_interactions:
            self.logger.warning(
                f"Number of pairs ({len(pairs)}) exceeds max_interactions ({self.max_interactions}). "
                "Consider using feature importance to select top interactions."
            )
            pairs = pairs[:self.max_interactions]

        for col1, col2 in pairs:
            feature_name = f'{col1}_{interaction_type}_{col2}'

            if interaction_type == 'multiply':
                df[feature_name] = df[col1] * df[col2]
            elif interaction_type == 'add':
                df[feature_name] = df[col1] + df[col2]
            elif interaction_type == 'subtract':
                df[feature_name] = df[col1] - df[col2]
            elif interaction_type == 'divide':
                # Avoid division by zero
                df[feature_name] = df[col1] / (df[col2] + 1e-8)
            else:
                raise ValueError(f"Unknown interaction type: {interaction_type}")

        return df

    def create_ratio_features(
        self,
        df: pd.DataFrame,
        numerator_columns: List[str],
        denominator_columns: List[str]
    ) -> pd.DataFrame:
        """
        Create ratio features.

        Args:
            df: Input DataFrame
            numerator_columns: Columns for numerator
            denominator_columns: Columns for denominator

        Returns:
            DataFrame with ratio features
        """
        df = df.copy()

        for num_col in numerator_columns:
            for den_col in denominator_columns:
                if num_col != den_col:
                    feature_name = f'{num_col}_div_{den_col}'
                    # Add small epsilon to avoid division by zero
                    df[feature_name] = df[num_col] / (df[den_col] + 1e-8)

        return df

    def create_polynomial_features(
        self,
        df: pd.DataFrame,
        numerical_columns: List[str]
    ) -> pd.DataFrame:
        """
        Create polynomial features using sklearn.

        Args:
            df: Input DataFrame
            numerical_columns: Columns to create polynomial features for

        Returns:
            DataFrame with polynomial features
        """
        if not self.include_polynomial:
            return df

        # Select numerical columns
        X = df[numerical_columns].values

        # Create polynomial features
        poly = PolynomialFeatures(
            degree=self.polynomial_degree,
            include_bias=self.include_bias,
            interaction_only=False
        )
        X_poly = poly.fit_transform(X)

        # Get feature names
        feature_names = poly.get_feature_names_out(numerical_columns)

        # Create DataFrame with polynomial features
        df_poly = pd.DataFrame(X_poly, columns=feature_names, index=df.index)

        # Remove original columns to avoid duplication
        original_cols = set(numerical_columns)
        new_cols = [col for col in df_poly.columns if col not in original_cols]
        df_poly = df_poly[new_cols]

        # Limit number of features
        if len(new_cols) > self.max_interactions:
            self.logger.warning(
                f"Polynomial features ({len(new_cols)}) exceed max_interactions. "
                "Keeping first {self.max_interactions} features."
            )
            df_poly = df_poly.iloc[:, :self.max_interactions]

        # Merge with original DataFrame
        df = pd.concat([df, df_poly], axis=1)

        return df

    def create_domain_specific_interactions(
        self,
        df: pd.DataFrame,
        interaction_rules: Dict[str, Tuple[List[str], Callable]]
    ) -> pd.DataFrame:
        """
        Create domain-specific interaction features.

        Args:
            df: Input DataFrame
            interaction_rules: Dictionary mapping feature names to (columns, function) tuples

        Returns:
            DataFrame with domain-specific interactions

        Example:
            interaction_rules = {
                'bmi': (['weight', 'height'], lambda w, h: w / (h ** 2)),
                'age_income_ratio': (['age', 'income'], lambda a, i: a / i)
            }
        """
        df = df.copy()

        for feature_name, (columns, func) in interaction_rules.items():
            # Extract column values
            column_values = [df[col].values for col in columns]

            # Apply function
            df[feature_name] = func(*column_values)

        return df

    def select_top_interactions(
        self,
        df: pd.DataFrame,
        target: pd.Series,
        numerical_columns: List[str],
        n_top: int = 50,
        method: str = 'mutual_info'  # 'mutual_info', 'correlation', 'random_forest'
    ) -> List[Tuple[str, ...]]:
        """
        Select top interaction features based on importance.

        Args:
            df: Input DataFrame
            target: Target variable
            numerical_columns: Numerical columns
            n_top: Number of top interactions to select
            method: Method for ranking interactions

        Returns:
            List of selected interaction pairs
        """
        from sklearn.feature_selection import mutual_info_regression, mutual_info_classif

        # Generate all pairwise combinations
        pairs = list(combinations(numerical_columns, 2))

        # Compute interaction importance
        importances = []

        for col1, col2 in pairs:
            interaction = df[col1] * df[col2]

            if method == 'mutual_info':
                # Determine if classification or regression
                if target.nunique() < 20:  # Heuristic for classification
                    mi = mutual_info_classif(interaction.values.reshape(-1, 1), target)[0]
                else:
                    mi = mutual_info_regression(interaction.values.reshape(-1, 1), target)[0]
                importances.append(mi)

            elif method == 'correlation':
                corr = np.abs(np.corrcoef(interaction, target)[0, 1])
                importances.append(corr if not np.isnan(corr) else 0)

            else:
                raise ValueError(f"Unknown method: {method}")

        # Select top pairs
        top_indices = np.argsort(importances)[::-1][:n_top]
        selected_pairs = [pairs[i] for i in top_indices]

        self.selected_interactions_ = selected_pairs

        self.logger.info(f"Selected {len(selected_pairs)} top interaction pairs")

        return selected_pairs

    def create_selected_interactions(
        self,
        df: pd.DataFrame,
        interaction_type: str = 'multiply'
    ) -> pd.DataFrame:
        """
        Create interactions from pre-selected pairs.

        Args:
            df: Input DataFrame
            interaction_type: Type of interaction

        Returns:
            DataFrame with selected interactions
        """
        if self.selected_interactions_ is None:
            raise ValueError("Run select_top_interactions first")

        df = df.copy()

        for col1, col2 in self.selected_interactions_:
            feature_name = f'{col1}_{interaction_type}_{col2}'

            if interaction_type == 'multiply':
                df[feature_name] = df[col1] * df[col2]
            elif interaction_type == 'add':
                df[feature_name] = df[col1] + df[col2]
            elif interaction_type == 'divide':
                df[feature_name] = df[col1] / (df[col2] + 1e-8)

        return df

# Example usage
if __name__ == "__main__":
    # Generate sample data
    np.random.seed(42)
    df = pd.DataFrame({
        'feature1': np.random.randn(1000),
        'feature2': np.random.randn(1000),
        'feature3': np.random.randn(1000),
        'feature4': np.random.randn(1000),
        'target': np.random.randn(1000)
    })

    numerical_cols = ['feature1', 'feature2', 'feature3', 'feature4']

    # Initialize generator
    generator = InteractionFeatureGenerator(
        max_interactions=100,
        include_polynomial=True,
        polynomial_degree=2
    )

    # Select top interactions
    selected_pairs = generator.select_top_interactions(
        df, df['target'], numerical_cols, n_top=10
    )

    print(f"\nTop 10 interaction pairs:")
    for pair in selected_pairs:
        print(f"  {pair[0]} × {pair[1]}")

    # Create selected interactions
    df_with_interactions = generator.create_selected_interactions(df)

    print(f"\nOriginal shape: {df.shape}")
    print(f"With interactions: {df_with_interactions.shape}")
\end{lstlisting}

\section{Categorical Encoding Strategies}

Proper categorical encoding is critical for model performance. Different encoding strategies are appropriate for different scenarios and model types.

\subsection{Advanced Categorical Encoding}

\begin{lstlisting}[style=python, caption={Comprehensive Categorical Encoding Framework}]
from typing import List, Dict, Optional, Union
import pandas as pd
import numpy as np
from sklearn.preprocessing import LabelEncoder, OneHotEncoder
import category_encoders as ce
import logging

class CategoricalEncoder:
    """
    Comprehensive categorical encoding framework.

    Encoding strategies:
    - Label Encoding
    - One-Hot Encoding
    - Target Encoding (mean encoding)
    - Frequency Encoding
    - Binary Encoding
    - Hash Encoding (for high cardinality)
    - Ordinal Encoding
    """

    def __init__(
        self,
        encoding_strategy: Dict[str, str] = None,
        handle_unknown: str = 'value',  # 'value', 'error', 'ignore'
        unknown_value: int = -1,
        smoothing: float = 1.0  # For target encoding
    ):
        """
        Initialize categorical encoder.

        Args:
            encoding_strategy: Dictionary mapping column names to encoding strategies
            handle_unknown: How to handle unknown categories
            unknown_value: Value for unknown categories
            smoothing: Smoothing parameter for target encoding
        """
        self.encoding_strategy = encoding_strategy or {}
        self.handle_unknown = handle_unknown
        self.unknown_value = unknown_value
        self.smoothing = smoothing

        # State
        self.encoders_: Dict[str, Any] = {}
        self.fitted_ = False

        self.logger = logging.getLogger(__name__)

    def fit(self, X: pd.DataFrame, y: Optional[pd.Series] = None) -> 'CategoricalEncoder':
        """
        Fit encoders on data.

        Args:
            X: Input DataFrame
            y: Target variable (required for target encoding)

        Returns:
            Self
        """
        for col, strategy in self.encoding_strategy.items():
            if col not in X.columns:
                self.logger.warning(f"Column {col} not found in DataFrame")
                continue

            if strategy == 'label':
                encoder = LabelEncoder()
                encoder.fit(X[col])
                self.encoders_[col] = encoder

            elif strategy == 'onehot':
                encoder = OneHotEncoder(
                    handle_unknown=self.handle_unknown,
                    sparse_output=False
                )
                encoder.fit(X[[col]])
                self.encoders_[col] = encoder

            elif strategy == 'target':
                if y is None:
                    raise ValueError("Target variable required for target encoding")
                encoder = ce.TargetEncoder(
                    cols=[col],
                    smoothing=self.smoothing,
                    handle_unknown=self.handle_unknown
                )
                encoder.fit(X[[col]], y)
                self.encoders_[col] = encoder

            elif strategy == 'frequency':
                # Frequency encoding: replace with frequency in dataset
                freq = X[col].value_counts(normalize=True).to_dict()
                self.encoders_[col] = freq

            elif strategy == 'binary':
                encoder = ce.BinaryEncoder(cols=[col])
                encoder.fit(X[[col]])
                self.encoders_[col] = encoder

            elif strategy == 'hash':
                # Hash encoding for high cardinality
                encoder = ce.HashingEncoder(
                    cols=[col],
                    n_components=8  # Number of hash bins
                )
                encoder.fit(X[[col]])
                self.encoders_[col] = encoder

            elif strategy == 'ordinal':
                encoder = ce.OrdinalEncoder(cols=[col])
                encoder.fit(X[[col]])
                self.encoders_[col] = encoder

            else:
                raise ValueError(f"Unknown encoding strategy: {strategy}")

        self.fitted_ = True
        return self

    def transform(self, X: pd.DataFrame) -> pd.DataFrame:
        """
        Transform data using fitted encoders.

        Args:
            X: Input DataFrame

        Returns:
            Transformed DataFrame
        """
        if not self.fitted_:
            raise ValueError("Encoder must be fitted before transform")

        X_transformed = X.copy()

        for col, encoder in self.encoders_.items():
            if col not in X_transformed.columns:
                continue

            strategy = self.encoding_strategy[col]

            if strategy == 'label':
                # Handle unknown categories
                X_transformed[col] = X_transformed[col].apply(
                    lambda x: encoder.transform([x])[0]
                    if x in encoder.classes_
                    else self.unknown_value
                )

            elif strategy == 'onehot':
                # One-hot encoding
                encoded = encoder.transform(X_transformed[[col]])
                feature_names = [f'{col}_{cat}' for cat in encoder.categories_[0]]

                # Remove original column and add encoded columns
                X_transformed = X_transformed.drop(columns=[col])
                for i, name in enumerate(feature_names):
                    X_transformed[name] = encoded[:, i]

            elif strategy == 'target':
                # Target encoding
                encoded = encoder.transform(X_transformed[[col]])
                X_transformed[col] = encoded[col].values

            elif strategy == 'frequency':
                # Frequency encoding
                freq_dict = encoder
                X_transformed[col] = X_transformed[col].map(freq_dict).fillna(0)

            elif strategy in ['binary', 'hash', 'ordinal']:
                # Category encoders
                encoded = encoder.transform(X_transformed[[col]])
                X_transformed = X_transformed.drop(columns=[col])
                X_transformed = pd.concat([X_transformed, encoded], axis=1)

        return X_transformed

    def fit_transform(self, X: pd.DataFrame, y: Optional[pd.Series] = None) -> pd.DataFrame:
        """Fit and transform in one step."""
        return self.fit(X, y).transform(X)

    def get_feature_names(self) -> List[str]:
        """Get names of all encoded features."""
        feature_names = []

        for col, encoder in self.encoders_.items():
            strategy = self.encoding_strategy[col]

            if strategy == 'onehot':
                for cat in encoder.categories_[0]:
                    feature_names.append(f'{col}_{cat}')
            elif strategy in ['binary', 'hash']:
                # These create multiple columns
                feature_names.extend([f'{col}_{i}' for i in range(encoder.n_features_out_)])
            else:
                # Single column encodings
                feature_names.append(col)

        return feature_names

# Example usage
if __name__ == "__main__":
    # Generate sample data
    np.random.seed(42)
    df = pd.DataFrame({
        'category_low_card': np.random.choice(['A', 'B', 'C'], 1000),
        'category_high_card': np.random.choice([f'cat_{i}' for i in range(100)], 1000),
        'category_ordinal': np.random.choice(['low', 'medium', 'high'], 1000),
        'target': np.random.randint(0, 2, 1000)
    })

    # Define encoding strategies
    encoding_config = {
        'category_low_card': 'onehot',
        'category_high_card': 'hash',
        'category_ordinal': 'target'
    }

    # Initialize encoder
    encoder = CategoricalEncoder(
        encoding_strategy=encoding_config,
        smoothing=1.0
    )

    # Fit and transform
    df_encoded = encoder.fit_transform(df, df['target'])

    print(f"\nOriginal shape: {df.shape}")
    print(f"Encoded shape: {df_encoded.shape}")
    print(f"\nOriginal columns: {df.columns.tolist()}")
    print(f"Encoded columns ({len(df_encoded.columns)} total)")
\end{lstlisting}

\section{Feature Selection and Importance}

Feature selection reduces dimensionality, improves model performance, and enhances interpretability by removing irrelevant or redundant features.

\subsection{Automated Feature Selection}

\begin{lstlisting}[style=python, caption={Comprehensive Feature Selection Framework}]
from typing import List, Optional, Dict, Any, Tuple
import pandas as pd
import numpy as np
from sklearn.feature_selection import (
    SelectKBest, mutual_info_classif, mutual_info_regression,
    f_classif, f_regression, RFE, SelectFromModel
)
from sklearn.ensemble import RandomForestClassifier, RandomForestRegressor
from sklearn.linear_model import LassoCV
import logging

class FeatureSelector:
    """
    Comprehensive feature selection framework.

    Methods:
    - Filter methods (univariate statistical tests)
    - Wrapper methods (RFE)
    - Embedded methods (L1 regularization, tree importance)
    - Correlation-based selection
    - Variance threshold
    """

    def __init__(
        self,
        task_type: str = 'classification',  # 'classification' or 'regression'
        selection_method: str = 'mutual_info',  # 'mutual_info', 'f_test', 'rfe', 'lasso', 'tree'
        n_features: Optional[int] = None,
        threshold: Optional[float] = None
    ):
        """
        Initialize feature selector.

        Args:
            task_type: Type of ML task
            selection_method: Feature selection method
            n_features: Number of features to select
            threshold: Threshold for feature importance
        """
        self.task_type = task_type
        self.selection_method = selection_method
        self.n_features = n_features
        self.threshold = threshold

        self.selector_ = None
        self.selected_features_: Optional[List[str]] = None
        self.feature_importance_: Optional[pd.DataFrame] = None

        self.logger = logging.getLogger(__name__)

    def fit(self, X: pd.DataFrame, y: pd.Series) -> 'FeatureSelector':
        """
        Fit feature selector.

        Args:
            X: Features
            y: Target

        Returns:
            Self
        """
        feature_names = X.columns.tolist()

        if self.selection_method == 'mutual_info':
            # Mutual information
            if self.task_type == 'classification':
                score_func = mutual_info_classif
            else:
                score_func = mutual_info_regression

            selector = SelectKBest(score_func=score_func, k=self.n_features or 'all')
            selector.fit(X, y)

            # Get scores
            scores = pd.DataFrame({
                'feature': feature_names,
                'score': selector.scores_
            }).sort_values('score', ascending=False)

            self.feature_importance_ = scores
            self.selector_ = selector

            # Select top features
            if self.n_features:
                self.selected_features_ = scores.head(self.n_features)['feature'].tolist()
            else:
                self.selected_features_ = feature_names

        elif self.selection_method == 'f_test':
            # F-test (ANOVA for classification, F-statistic for regression)
            if self.task_type == 'classification':
                score_func = f_classif
            else:
                score_func = f_regression

            selector = SelectKBest(score_func=score_func, k=self.n_features or 'all')
            selector.fit(X, y)

            scores = pd.DataFrame({
                'feature': feature_names,
                'score': selector.scores_
            }).sort_values('score', ascending=False)

            self.feature_importance_ = scores
            self.selector_ = selector
            self.selected_features_ = scores.head(self.n_features or len(feature_names))['feature'].tolist()

        elif self.selection_method == 'rfe':
            # Recursive Feature Elimination
            if self.task_type == 'classification':
                estimator = RandomForestClassifier(n_estimators=50, random_state=42)
            else:
                estimator = RandomForestRegressor(n_estimators=50, random_state=42)

            selector = RFE(
                estimator=estimator,
                n_features_to_select=self.n_features,
                step=1
            )
            selector.fit(X, y)

            # Get feature ranking
            rankings = pd.DataFrame({
                'feature': feature_names,
                'ranking': selector.ranking_
            }).sort_values('ranking')

            self.feature_importance_ = rankings
            self.selector_ = selector
            self.selected_features_ = [f for f, r in zip(feature_names, selector.support_) if r]

        elif self.selection_method == 'lasso':
            # L1 regularization (Lasso)
            lasso = LassoCV(cv=5, random_state=42, max_iter=10000)
            lasso.fit(X, y)

            # Get coefficients
            coefs = pd.DataFrame({
                'feature': feature_names,
                'coefficient': np.abs(lasso.coef_)
            }).sort_values('coefficient', ascending=False)

            self.feature_importance_ = coefs
            self.selector_ = lasso

            # Select features with non-zero coefficients
            if self.threshold:
                self.selected_features_ = coefs[coefs['coefficient'] > self.threshold]['feature'].tolist()
            elif self.n_features:
                self.selected_features_ = coefs.head(self.n_features)['feature'].tolist()
            else:
                self.selected_features_ = coefs[coefs['coefficient'] > 0]['feature'].tolist()

        elif self.selection_method == 'tree':
            # Tree-based feature importance
            if self.task_type == 'classification':
                model = RandomForestClassifier(n_estimators=100, random_state=42)
            else:
                model = RandomForestRegressor(n_estimators=100, random_state=42)

            model.fit(X, y)

            importances = pd.DataFrame({
                'feature': feature_names,
                'importance': model.feature_importances_
            }).sort_values('importance', ascending=False)

            self.feature_importance_ = importances
            self.selector_ = model

            if self.n_features:
                self.selected_features_ = importances.head(self.n_features)['feature'].tolist()
            elif self.threshold:
                self.selected_features_ = importances[
                    importances['importance'] > self.threshold
                ]['feature'].tolist()
            else:
                self.selected_features_ = feature_names

        else:
            raise ValueError(f"Unknown selection method: {self.selection_method}")

        self.logger.info(f"Selected {len(self.selected_features_)} features using {self.selection_method}")

        return self

    def transform(self, X: pd.DataFrame) -> pd.DataFrame:
        """
        Transform data by selecting features.

        Args:
            X: Features

        Returns:
            DataFrame with selected features
        """
        if self.selected_features_ is None:
            raise ValueError("Selector must be fitted before transform")

        return X[self.selected_features_]

    def fit_transform(self, X: pd.DataFrame, y: pd.Series) -> pd.DataFrame:
        """Fit and transform in one step."""
        return self.fit(X, y).transform(X)

    def get_feature_importance(self) -> pd.DataFrame:
        """Get feature importance scores."""
        if self.feature_importance_ is None:
            raise ValueError("Selector must be fitted first")

        return self.feature_importance_

    def remove_correlated_features(
        self,
        X: pd.DataFrame,
        threshold: float = 0.95
    ) -> List[str]:
        """
        Remove highly correlated features.

        Args:
            X: Features
            threshold: Correlation threshold

        Returns:
            List of features to keep
        """
        # Compute correlation matrix
        corr_matrix = X.corr().abs()

        # Find pairs of highly correlated features
        upper_tri = corr_matrix.where(
            np.triu(np.ones(corr_matrix.shape), k=1).astype(bool)
        )

        # Find features to drop
        to_drop = [
            column for column in upper_tri.columns
            if any(upper_tri[column] > threshold)
        ]

        features_to_keep = [f for f in X.columns if f not in to_drop]

        self.logger.info(f"Removed {len(to_drop)} correlated features (threshold={threshold})")

        return features_to_keep

    def remove_low_variance_features(
        self,
        X: pd.DataFrame,
        threshold: float = 0.01
    ) -> List[str]:
        """
        Remove features with low variance.

        Args:
            X: Features
            threshold: Variance threshold

        Returns:
            List of features to keep
        """
        from sklearn.feature_selection import VarianceThreshold

        selector = VarianceThreshold(threshold=threshold)
        selector.fit(X)

        features_to_keep = X.columns[selector.get_support()].tolist()

        self.logger.info(
            f"Removed {X.shape[1] - len(features_to_keep)} low variance features"
        )

        return features_to_keep

    def plot_feature_importance(self, top_n: int = 20, save_path: Optional[str] = None):
        """Plot feature importance."""
        if self.feature_importance_ is None:
            raise ValueError("No feature importance available")

        import matplotlib.pyplot as plt

        top_features = self.feature_importance_.head(top_n)

        plt.figure(figsize=(10, 8))
        plt.barh(range(len(top_features)), top_features.iloc[:, 1].values)
        plt.yticks(range(len(top_features)), top_features.iloc[:, 0].values)
        plt.xlabel('Importance')
        plt.title(f'Top {top_n} Features ({self.selection_method})')
        plt.gca().invert_yaxis()
        plt.tight_layout()

        if save_path:
            plt.savefig(save_path, dpi=150, bbox_inches='tight')
        else:
            plt.savefig('feature_importance.png', dpi=150, bbox_inches='tight')
        plt.close()

# Example usage
if __name__ == "__main__":
    from sklearn.datasets import make_classification

    # Generate sample data
    X, y = make_classification(
        n_samples=1000,
        n_features=50,
        n_informative=20,
        n_redundant=10,
        random_state=42
    )
    X = pd.DataFrame(X, columns=[f'feature_{i}' for i in range(50)])
    y = pd.Series(y)

    # Test different selection methods
    methods = ['mutual_info', 'f_test', 'rfe', 'tree']

    for method in methods:
        print(f"\n=== {method.upper()} ===")

        selector = FeatureSelector(
            task_type='classification',
            selection_method=method,
            n_features=20
        )

        X_selected = selector.fit_transform(X, y)

        print(f"Selected {X_selected.shape[1]} features")
        print(f"\nTop 10 features:")
        print(selector.get_feature_importance().head(10))

        # Plot importance
        selector.plot_feature_importance(top_n=15, save_path=f'importance_{method}.png')

    # Remove correlated features
    selector = FeatureSelector(task_type='classification')
    features_to_keep = selector.remove_correlated_features(X, threshold=0.9)
    print(f"\n{len(features_to_keep)} features after removing correlated (threshold=0.9)")
\end{lstlisting}

\section{Chapter Summary}

This chapter presented comprehensive approaches to feature engineering at scale:

\begin{itemize}
    \item \textbf{Feature Engineering Framework:} Modular pipeline with versioning, metadata tracking, and train/test consistency

    \item \textbf{Temporal Features:} Basic components, cyclical encodings, business calendar features, and relative time calculations

    \item \textbf{Aggregation Features:} Group-by aggregations, rolling windows, expanding windows, lag features, and time-based aggregations

    \item \textbf{Interaction Features:} Pairwise interactions, polynomial features, ratio features, and intelligent selection

    \item \textbf{Categorical Encoding:} Label encoding, one-hot encoding, target encoding, frequency encoding, and hash encoding for high cardinality

    \item \textbf{Feature Selection:} Filter methods, wrapper methods, embedded methods, correlation analysis, and variance thresholding
\end{itemize}

\subsection{Key Takeaways}

\begin{enumerate}
    \item \textbf{Systematic Approach:} Use structured frameworks for feature engineering to ensure reproducibility and maintainability

    \item \textbf{Train-Test Consistency:} Always fit encoders and transformations on training data and apply to test/serving data

    \item \textbf{Domain Knowledge:} Leverage business understanding to create meaningful domain-specific features

    \item \textbf{Avoid Data Leakage:} Never use future information or target-derived features in ways that wouldn't be available at prediction time

    \item \textbf{Manage Dimensionality:} Use feature selection to reduce curse of dimensionality and improve model performance

    \item \textbf{Handle Missing Values:} Explicit strategy for missing data with indicator features when appropriate

    \item \textbf{Monitor Features:} Track feature distributions and drift in production to detect data quality issues

    \item \textbf{Optimize Computation:} Consider computational cost for real-time serving, especially for complex aggregations
\end{enumerate}

\section{Further Reading}

\begin{itemize}
    \item Kuhn, M., \& Johnson, K. (2019). \textit{Feature Engineering and Selection: A Practical Approach for Predictive Models}
    \item Zheng, A., \& Casari, A. (2018). \textit{Feature Engineering for Machine Learning}
    \item Cerda, P., \& Varoquaux, G. (2020). Encoding high-cardinality string categorical variables. \textit{IEEE TKDE}
    \item Micci-Barreca, D. (2001). A preprocessing scheme for high-cardinality categorical attributes. \textit{SIGKDD Explorations}
    \item Guyon, I., \& Elisseeff, A. (2003). An introduction to variable and feature selection. \textit{JMLR}
\end{itemize}

\section{Exercises}

\begin{enumerate}
    \item \textbf{Temporal Feature Engineering:}
    \begin{itemize}
        \item Create a comprehensive temporal feature set for a time series dataset
        \item Include cyclical encodings for multiple time periods
        \item Add business calendar features and holidays
        \item Validate that cyclical features correctly handle period boundaries
    \end{itemize}

    \item \textbf{Aggregation Features:}
    \begin{itemize}
        \item Implement rolling window features with multiple window sizes
        \item Create lag features at different time horizons
        \item Compute expanding window statistics
        \item Ensure proper handling of group-wise aggregations
    \end{itemize}

    \item \textbf{Feature Interactions:}
    \begin{itemize}
        \item Generate pairwise interaction features
        \item Use feature importance to select top interactions
        \item Create domain-specific interaction features based on problem knowledge
        \item Manage dimensionality to avoid explosion
    \end{itemize}

    \item \textbf{Categorical Encoding:}
    \begin{itemize}
        \item Implement multiple encoding strategies for different categorical variables
        \item Compare target encoding with and without smoothing
        \item Use hash encoding for high-cardinality features
        \item Evaluate impact of different encodings on model performance
    \end{itemize}

    \item \textbf{Feature Selection:}
    \begin{itemize}
        \item Apply multiple feature selection methods (filter, wrapper, embedded)
        \item Compare selected features across different methods
        \item Remove highly correlated features
        \item Evaluate model performance vs. number of features
    \end{itemize}

    \item \textbf{End-to-End Feature Engineering:}
    \begin{itemize}
        \item Build a complete feature engineering pipeline from raw data to model-ready features
        \item Include temporal, aggregation, interaction, and encoding features
        \item Apply feature selection to reduce dimensionality
        \item Ensure train/test consistency throughout
        \item Track feature metadata and provenance
    \end{itemize}

    \item \textbf{Production Feature Pipeline:}
    \begin{itemize}
        \item Design a feature engineering pipeline for online serving
        \item Optimize computation for low-latency requirements
        \item Implement feature versioning and rollback capability
        \item Add monitoring for feature drift
        \item Document all features with metadata
    \end{itemize}
\end{enumerate}
